% from ega3-1-fr.tex
\oldpage[III]{115}
\begin{env}[2.6.3]
\label{III.2.6.3-fr}
\emph{Remarque.} Lorsque $X$ est propre sur $Y$, on peut démontrer
(\hyperref[III.2.6.2-fr]{2.6.2}) plus simplement, en raisonnant comme dans
le th. de Chevalley (II, 6.7.1), à l'aide de
(\hyperref[III.2.6.1-fr]{2.6.1}) et du lemme (II, 6.7.1.1).
\end{env}

\section{Le théorème de finitude pour les morphismes propres}
\label{section:III.3-fr}

\subsection{Le lemme de dévissage}
\label{subsection:III.3.1-fr}

\begin{definition}[3.1.1]
\label{III.3.1.1-fr}
Soit $K$ une catégorie abélienne. On dit qu'un sous-ensemble $K'$ de
l'ensemble des objets de $K$ est \emph{exact} si $0\in K'$ et si, pour toute
suite exacte
\[
0\longrightarrow A'\longrightarrow A\longrightarrow A''\longrightarrow0
\]
dans $K$ telle que deux des objets $A$, $A'$, $A''$ soit dans $K'$, alors le
troisième est aussi dans $K'$.
\end{definition}

\begin{theorem}[3.1.2]
\label{III.3.1.2-fr}
Soit $X$ un préschéma noethérien~; on désigne par $K$ la catégorie abélienne
des $\mathcal O_X$-Modules cohérents. Soient $K'$ un sous-ensemble exact de
$K$, $X'$ une partie fermée de l'espace sous-jacent à $X$. On suppose que
pour toute partie fermée irréductible $Y$ de $X'$, de point générique $y$, il
existe un $\mathcal O_X$-Module $\mathcal G\in K'$ tel que $\mathcal G_y$ soit
un $k(y)$-espace vectoriel de dimension $1$. Alors tout $\mathcal O_X$-Module
cohérent de support contenu dans $X'$ appartient à $K'$ (et en particulier,
si $X'=X$, on a $K'=K$).
\end{theorem}

\begin{proof}
Considérons la propriété suivante $P(Y)$ d'une partie fermée $Y$ de $X'$~:
tout $\mathcal O_X$-Module cohérent de support contenu dans $Y$ appartient à
$K'$. En vertu du principe de récurrence noethérienne (0\textsubscript{I},
2.2.2), on voit qu'on est ramené à démontrer que \emph{si $Y$ est une partie
fermée de $X'$ telle que la propriété $P(Y')$ soit vraie pour toute partie
fermée $Y'$ de $Y$, distincte de $Y$, alors $P(Y)$ est vraie}.

Soit donc $\mathcal F\in K$ à support contenu dans $Y$ et prouvons que
$\mathcal F\in K'$. Désignons encore par $Y$ le sous-préschéma fermé réduit de
$X$ ayant $Y$ pour espace sous-jacent (I, 5.2.1)~; il est défini par un Idéal
cohérent $\mathcal J$ de $\mathcal O_X$. On sait (I, 9.3.4) qu'il existe un
entier $n>0$ tel que $\mathcal J^n\mathcal F=0$~; pour $1\leq k\leq n$, on a
donc une suite exacte
\[
0\longrightarrow
\mathcal J^{k-1}\mathcal F/\mathcal J^k\mathcal F
\longrightarrow \mathcal F/\mathcal J^k\mathcal F
\longrightarrow \mathcal F/\mathcal J^{k-1}\mathcal F
\longrightarrow0
\]
de $\mathcal O_X$-Modules cohérents (0\textsubscript{I}, 5.3.6 et 5.3.3)~;
comme $K'$ est exact, on voit, par récurrence sur $k$, qu'il suffit de montrer
que chacun des
$\mathcal F_k=\mathcal J^{k-1}\mathcal F/\mathcal J^k\mathcal F$ est dans
$K'$. On est donc ramené à prouver que $\mathcal F\in K'$ sous l'hypothèse
supplémentaire que $\mathcal J\mathcal F=0$~; il revient au même de dire que
$\mathcal F=j_*(j^*(\mathcal F))$, où $j$ est l'injection canonique
$Y\to X$. Distinguons maintenant deux cas~:
\begin{enumerate}
\item[a)] $Y$ est \emph{réductible}. Soient $Y=Y'\cup Y''$, $Y'$ et $Y''$
étant des parties fermées de $Y$, distinctes de $Y$~; soient encore $Y'$,
$Y''$ les sous-préschémas fermés réduits de $X$ ayant pour espaces
sous-jacents $Y'$, $Y''$ respectivement, qui sont définis respectivement par
les Idéaux cohérents $\mathcal J'$, $\mathcal J''$ de $\mathcal O_X$. Posons
\[
\mathcal F'=\mathcal F\otimes_{\mathcal O_X}
(\mathcal O_X/\mathcal J'),\qquad
\mathcal F''=\mathcal F\otimes_{\mathcal O_X}
(\mathcal O_X/\mathcal J'').
\]
Les homomorphismes canoniques $\mathcal F\to\mathcal F'$,
$\mathcal F\to\mathcal F''$ définissent donc un homomorphisme
$u:\mathcal F\to\mathcal F'\oplus\mathcal F''$. Montrons que pour tout
$z\notin Y'\cap Y''$, l'homomorphisme
$u_z:\mathcal F_z\to\mathcal F'_z\oplus\mathcal F''_z$ est \emph{bijectif}.
En effet, on a $\mathcal J'\cap\mathcal J''=\mathcal J$, car la question est
locale et

